% =========================================================================
% ELEMENTOS TEXTUAIS E METODOLÓGICOS DA INICIAÇÃO CIENTÍFICA
% =========================================================================
\section{Introdução}\label{sec:introducao}
A Otimização Combinatória é uma área de pesquisa situada na interseção entre Matemática e Ciência da Computação (veja~\cite{korte2018combinatorial}). Do ponto de vista matemático, a área mantém forte influência com a Teoria dos Grafos, uma vez que digrafos constituem objetos combinatórios com amplo poder de modelagem (veja~\cite{bondy2008graph, diestel2017graph}).
Sob a perspectiva computacional, o foco reside na obtenção de algoritmos eficientes que produzam soluções garantidamente ótimas (veja~\cite{schrijver2003combinatorial, cook1998combinatorial}). Dentre os tópicos de destaque, os problemas de fluxos em digrafos ocupam posição central, tratando de redes onde recursos devem ser transportados respeitando restrições de capacidade (veja~\cite{ahuja1993network}).
O objetivo deste projeto é proporcionar uma formação sólida na área, integrando teoria e prática, desde as estruturas fundamentais de fluxo máximo até os modelos avançados de fluxo de custo mínimo e suas validações experimentais.

\subsection{Estrutura do Documento e Guia de Leitura para Avaliação}\label{sec:leitura}

Com a finalidade de facilitar a leitura crítica e proporcionar um roteiro claro de navegação aos avaliadores deste relatório, detalhamos a organização macroscópica do volume:

\begin{enumerate}[leftmargin=1.5em, itemsep=0.15em, topsep=0.2em, parsep=0em]
    \item \textbf{Corpo Analítico Principal (Seções~\ref{sec:introducao} a~\ref{sec:conclusoes}, páginas~\pageref{sec:introducao} a~\pageref{sec:conclusoes}):} Constrói o percurso teórico-matemático, formaliza as demonstrações de todos os teoremas estruturais fundamentais, detalha o funcionamento dos algoritmos de fluxo máximo e custo mínimo, apresenta as reduções combinatórias, a modelagem de segmentação de imagens e a metodologia experimental.
\ifapendicecodigo
    \item \textbf{Apêndices Computacionais (Apêndices~\ref{ap:flownetwork} a~\ref{ap:github}, páginas~\pageref{ap:flownetwork} a~\pageref{ap:github}):} Documentam o código-fonte integral desenvolvido em linguagem C++23, evidenciando o encapsulamento orientado a objetos, a robustez das estruturas de dados e a implementação dos motores algorítmicos.
\else
    \item \textbf{Repositório Aberto e Apêndices de Código:} O código-fonte integral em C++23, instâncias DIMACS e scripts de teste encontram-se documentados na versão estendida deste relatório, na monografia integral (\texttt{book.pdf}) e versionados no repositório aberto do projeto no GitHub.
\fi
\end{enumerate}

\noindent \textbf{Roteiro de Leitura Rápida (\textit{Fast Track} para o Avaliador):}
Caso o parecerista disponha de tempo reduzido para a análise e deseje examinar pontualmente as contribuições mais expressivas do projeto, recomenda-se a consulta prioritária aos seguintes núcleos:
\begin{itemize}[leftmargin=1.5em, itemsep=0.15em, topsep=0.2em, parsep=0em]
    \item \textbf{Subseção~\ref{sec:push_relabel_gap} (página~\pageref{sec:push_relabel_gap}):} Aceleração prática de convergência por meio da heurística de lacuna (\textit{Gap Heuristic}) no algoritmo Push-Relabel;
    \item \textbf{Subseção~\ref{sec:segmentacao} (página~\pageref{sec:segmentacao}):} Modelagem e implementação prática de Segmentação Interativa de Imagens via corte mínimo em grafos;
    \item \textbf{Subseção~\ref{sec:ssp_dijkstra} (página~\pageref{sec:ssp_dijkstra}):} Paradigma Dual de Custo Mínimo via \textit{Successive Shortest Path}, potenciais nodais de Edmonds-Karp e invariante de não-negatividade ($\bar{w}_\pi \ge 0$) com Dijkstra;
    \item \textbf{Subseção~\ref{sec:network_simplex} (página~\pageref{sec:network_simplex}):} Teoria da Total Unimodularidade e mecânica computacional do algoritmo \textit{Network Simplex} sobre árvores geradoras;
    \item \textbf{Subseção~\ref{sec:benchmarks} (página~\pageref{sec:benchmarks}):} Metodologia de validação empírica e protocolo experimental em instâncias padronizadas DIMACS;
    \item \textbf{Seção~\ref{sec:conclusoes} (página~\pageref{sec:conclusoes}):} Síntese conclusiva das metas alcançadas e das diretrizes de pesquisa futura.
\end{itemize}

\section{Fundamentação teórica}\label{sec:fundamentacao}
O problema do $st$-fluxo máximo, onde se busca maximizar a quantidade de fluxo enviada de um vértice fonte $s$ para um vértice de destino $t$, é um dos pilares deste estudo (veja~\cite{ford1956maximal}). Foram analisados algoritmos como os de Ford-Fulkerson, Edmonds-Karp, Dinic e Push-Relabel (Goldberg-Tarjan) (veja~\cite{ford1956maximal, edmonds1972theoretical, dinic1970algorithm, goldberg1988new}).
O projeto também contempla o problema de fluxo de custo mínimo, que busca minimizar o custo total enquanto satisfaz as demandas da rede. Para este fim, foram desenvolvidos e analisados algoritmos clássicos como Cycle Canceling e Successive Shortest Path (com SPFA e Dijkstra acelerado por potenciais nodais), culminando na implementação do algoritmo \textit{network simplex}, uma especialização do algoritmo simplex sobre grafos (veja~\cite{dantzig1951application, ahuja1993network}).
A fundamentação teórica revisada nesta pesquisa reafirma que muitos problemas combinatórios, como o emparelhamento máximo em digrafos bipartidos, podem ser resolvidos via reduções para problemas de fluxo (veja~\cite{ahuja1993network}).

\section{Metodologia}\label{sec:metodologia}
\subsection{Materiais e Métodos}
A metodologia baseia-se em ciclos de estudo bibliográfico seguidos de implementação. Os algoritmos foram integralmente desenvolvidos em linguagem C++23 orientada a objetos para garantir alta performance, utilizando listas de adjacência e representações residuais otimizadas.
O protocolo de validação experimental foi estruturado sobre instâncias padronizadas das coleções DIMACS (veja~\cite{Jensen2022}), permitindo a comparação sistemática de desempenho. A análise fundamenta-se em métricas de tempo de execução de CPU e corretude dos algoritmos implementados frente às soluções ótimas conhecidas (veja~\cite{ford1956maximal}).

\subsection{Etapas da pesquisa}
As atividades do projeto foram estruturadas e encontram-se no seguinte status:
\begin{itemize}
    \item \textbf{Atividade 1:} Estudo dos fundamentos de Otimização Combinatória e Teoria dos Grafos (veja~\cite{cook1998combinatorial}). \textbf{[Concluída]}
    \item \textbf{Atividade 2:} Estudo teórico do problema de fluxo máximo e algoritmos clássicos (veja~\cite{ford1956maximal, edmonds1972theoretical, dinic1970algorithm, goldberg1988new}). \textbf{[Concluída]}
    \item \textbf{Atividade 3:} Implementação em C++ dos algoritmos de fluxo máximo. \textbf{[Concluída]}
    \item \textbf{Atividade 4:} Estudo de técnicas de modelagem e redução de problemas. \textbf{[Concluída]}
    \item \textbf{Atividade 5:} Implementação de problemas modelados (Emparelhamento e Caminhos Disjuntos). \textbf{[Concluída]}
    \item \textbf{Atividade 6:} Elaboração do relatório final (este documento). \textbf{[Concluída]}
    \item \textbf{Atividade 7:} Estudo teórico do problema de fluxo de custo mínimo (veja~\cite{cook1998combinatorial}). \textbf{[Concluída]}
    \item \textbf{Atividade 8:} Implementação em C++ do algoritmo \textit{network simplex}. \textbf{[Concluída]}
    \item \textbf{Atividade 9:} Testes experimentais em instâncias de larga escala (DIMACS). \textbf{[Estruturada]}
    \item \textbf{Atividade 10:} Análise comparativa de desempenho dos algoritmos. \textbf{[Estruturada]}
    \item \textbf{Atividade 11:} Elaboração do relatório final e consolidação dos resultados. \textbf{[Concluída]}
\end{itemize}

O encadeamento estrutural deste plano de trabalho é sumarizado no pipeline metodológico da Figura~\ref{fig:pipeline_metodologico}.

\begin{figure}[!ht]
\centering
\resizebox{\textwidth}{!}{%
\begin{tikzpicture}[thick, node distance=0.6cm]
  \node[stage box] (s1) at (0,0) {\textbf{1. Teoria \& Dualidade}\\Max-Flow Min-Cut\\Reduções Bipartidas\\Fluxo Custo Mínimo};
  \node[stage box] (s2) at (3.3,0) {\textbf{2. Arquitetura C++23}\\Herança \& Polimorfismo\\\texttt{FlowNetwork}\\\texttt{CostNetwork}};
  \node[stage box] (s3) at (6.6,0) {\textbf{3. Motores Algorítmicos}\\Edmonds-Karp, Dinic\\Push-Relabel (Heur.)\\SSP, Network Simplex};
  \node[stage box] (s4) at (9.9,0) {\textbf{4. Certificação Cruzada}\\Invariantes de Fluxo\\Dualidade Forte\\Google Test};
  \node[stage box] (s5) at (13.2,0) {\textbf{5. Benchmarks}\\Grades 2D, Camadas\\Grafos Aleatórios\\Segmentação Imagens};

  \draw[->, line width=1.5pt, blue!80!black] (s1) -- (s2);
  \draw[->, line width=1.5pt, blue!80!black] (s2) -- (s3);
  \draw[->, line width=1.5pt, blue!80!black] (s3) -- (s4);
  \draw[->, line width=1.5pt, blue!80!black] (s4) -- (s5);
\end{tikzpicture}%
}
\caption{Pipeline metodológico da pesquisa: fluxo integrado conectando a fundamentação teórica em otimização, a modelagem primal-dual, a arquitetura de software orientada a objetos em C++23, os procedimentos de certificação cruzada de corretude e os testes de desempenho computacional.}
\label{fig:pipeline_metodologico}
\end{figure}
